\section{Discussion}\label{sec:discussion}

\subsection{Risk, not prediction}

The clearest pattern in the results is that the system's value came from the risk side. The direction forecasts were no better than chance, yet the Risk Engine cut drawdown by about three quarters in both test periods. These two findings fit together once you look at what the Risk Engine actually does. With a direction signal close to random, every trade is close to a fair bet minus costs. The best that risk rules can do with fair bets is make the bad ones small, and that is what happened. The engine stayed out on days when the forecast was weak or the jump risk was high, so it traded on about a third of days. When it did trade, it took smaller positions and gave the stop more room on days it expected to be volatile, and larger positions with tighter stops on calm days. The fixed rules traded every day at the same size, so they took the full force of every bad week.

This is also why volatility could be forecast when direction could not. Large moves cluster \cite{cont2001}, so the size of recent moves says something real about the size of tomorrow's, even though it says nothing about the sign. The same logic is behind volatility-managed portfolios \cite{moreira2017}. The Beta Agent's forecast was weak, with a correlation of 0.19 and a clear upward bias, but it did not need to be strong. The tuned constants $k$ and $c$ absorb a constant bias, and the rules only need the forecast to rank days roughly correctly. A rescaled GARCH(1,1) forecast was more accurate than the Beta Agent, so the deep learning part of the risk side was not where the value came from. The value came from using a volatility forecast at all. Whether the Risk Engine would do even better with GARCH as its input was not tested, and is the most obvious next experiment. The rule ablation backs this up, since the volatility stop and volatility sizing were two of the three most useful rules.

There is a less flattering way to describe the same result. Part of the drawdown cut comes simply from being out of the market more often, and a system that never trades has no drawdown at all. The ablation shows that sitting out is not the whole story. Removing the weak-prediction filter almost doubled the number of trades, from 63 to 121, but only added 3.3 points of drawdown, while making stops and sizing fixed cost more. Still, the Risk Engine's returns were small in absolute terms, 4.5\% in 2025 and 2.4\% in 2026 at 2x leverage. What it showed is that losses can be controlled, not that the system makes money.

\subsection{Why the forecasting hypotheses were weak}

H1 and H2 assumed there was a directional signal for extra features and attention to find. The direction results suggest there was very little. Daily Bitcoin returns behaved close to a random walk throughout, consistent with the efficient market view that public information is already in the price \cite{fama1970}. When the signal is that weak, a larger model mostly has more room to fit noise. The training set was also small for a deep model, at about 1{,}500 trading days, and the two-layer Bi-LSTM trained badly on several seeds. Attention mainly helped by making training more stable, which is a real benefit but not the one H2 predicted. Training on MSE made things worse, because on a near-noise target the safest MSE prediction is close to zero, which left some seeds predicting the same sign every day.

The volatility target problem is worth stating as a general lesson. A forecast should always be compared with persistence before its RMSE is trusted. The 14-day target looked reasonable in the proposal and gave a believable-looking RMSE, and only the comparison with the naive forecast showed that the model was far worse than doing nothing. Any target built from overlapping windows has this problem. The GARCH comparison taught a second version of the same lesson. My first comparison scored GARCH's standard deviation against the absolute return, and made the Beta Agent look better than it was. Two forecasts are only comparable if they are trying to predict the same quantity.

\subsection{The liquidation result}

H3 predicted fewer liquidations, and the stress test showed the opposite at 20x. The cause is a design flaw in rule 4. A stop placed at $c\,\hat\sigma_{t+1}$ grows on volatile days, and at high leverage it can end up further from the entry than the liquidation level, which sits at about $1/L$. At 20x that is 5\%, while the volatility stop on a wild day could be wider. The broker then closes the trade first and takes the whole margin. The fix is simple: cap the stop distance below the liquidation distance, for example at 80\% of $1/L$. I did not make this change, because it was found on the test year, and changing the engine after seeing test results would amount to tuning on the test set. It is left as future work.

At realistic retail leverage the question did not really arise. Neither engine was liquidated at 2x, 5x or 10x, because any stop tighter than $1/L$ fires first and Bitcoin rarely gaps through stops. The backtest fills stops at the stop price without slippage, so in live trading stops would sometimes fill worse, and liquidation risk is somewhat understated here for both engines.

\subsection{Deviations from the proposal}\label{sec:deviations}

My supervisor confirmed that changes from the proposal are acceptable if they are explained, and the main ones are listed here. The volatility target changed from 14-day realised volatility to the size of the next day's move, for the reasons in \cref{sec:vol-target}. The Beta Agent became two separate GRUs instead of one GRU with two heads, because the heads interfered with each other. Rule 2 measures forecast strength as a percentile of the model's recent forecasts instead of a fixed size, because the seeds forecast on different scales. The data starts on 12 February 2018 rather than 1 January, because of the warm-up period. The hypotheses now refer to the log-return, as the models always did, rather than the closing price. The GARCH comparison was corrected late in the project, after I realised the original version compared different quantities (\cref{sec:res-vol}). Beyond the proposal, I added the 2026 forward test, the block bootstrap, the rule ablation, the leverage stress test, the forecaster swap and the direction baselines.

The strength filter and the decision to add a forward test were both made after looking at 2025 results. This is why I treat 2026 as the cleanest test in the project. It did not influence any decision, and the drawdown result held on it.

\subsection{Limitations}

The study covers one asset over one test year and nine months of forward data, so all of the conclusions are about Bitcoin over 2025 and 2026. Results varied a lot between seeds, with a standard deviation of about 1 in the Sharpe ratio, so any single run can mislead. The bootstrap captures uncertainty about which market days happened to occur, but not uncertainty from training or tuning, since the settings were fixed. The backtest does not model slippage or the overnight financing charges that CFD brokers apply to positions held overnight, and both would reduce returns. EU rules also require brokers to close a retail client's positions once the account falls to half of the initial margin \cite{esma2018}, while the backtest only closes a trade when the whole margin is gone, so it is slightly more lenient than a real EU broker. The rand balances ignore exchange-rate moves, and a R1{,}000 account would fall below some brokers' minimum trade size at these settings. Finally, the jump forecasts were poorly calibrated and the volatility forecasts biased upwards, so neither should be read as literal probabilities or volatility levels.
